\documentclass[level=1]{kidscode}
\begin{document}
\kctitle{Notes for Families}{What we are doing, how we keep children safe, and how you can help at home}

\section{Dear families,}
Welcome to our home coding club! Once a week, for one hour, a small group of up to four children
(about 10--11 years old) will learn to code together at our home. No experience is needed: we start
from zero, and everyone learns at their own pace.

The course has three levels of 10--12 weeks each, with one session a week. \textbf{Level 1 (this one) is called
\enquote{Introduction}}:
\begin{itemize}[itemsep=2pt]
  \item \textbf{Weeks 1--6: Scratch.} Children snap colored blocks together to make sprites move, draw,
        react and keep score. We use the free Scratch app on the computer, \textbf{offline}.
  \item \textbf{Weeks 7--12: Python.} Children move from blocks to typing real code with Python's
        \enquote{turtle}, which draws pictures on the screen. We use a free program called Thonny.
  \item \textbf{Weeks 10--12: a final project}, \enquote{My Interactive Art Show}: a program that asks
        the viewer questions and draws a picture that changes with their answers.
  \item \textbf{Week 12: a showcase} where the children present their projects to families. You are invited!
\end{itemize}
Every session has the same shape: a warm-up game, a new idea, a short demonstration, building
together in pairs, and a wrap-up where children show their work. Many sessions include an
\enquote{unplugged} activity without computers (games, cards, acting out code), because the ideas
behind coding are as important as typing. The children also learn to work as a pair: one
\enquote{drives} the keyboard and mouse, the other \enquote{navigates}, and they swap.

\section{Our safety and privacy promises}
\begin{safety}[What we promise]
\begin{itemize}[itemsep=2pt]
  \item \textbf{No accounts, no logins, no email addresses.} We use the offline Scratch app and Thonny.
        Nobody signs up for anything.
  \item \textbf{No personal information in projects.} Children use a nickname in their work, never their
        real full name, address, school, age or photos. Programs never ask for personal information either.
  \item \textbf{Nothing is uploaded or shared online.} Projects stay on the computers in the room. At the
        showcase, projects are shown on the screen only.
  \item \textbf{Photos only with permission.} We do not take photos of children. If you take photos at
        the showcase, please keep them in the family, and only with the child's agreement.
  \item \textbf{Healthy screen time.} Every session includes time away from the screen, a stretch break,
        and the 20-20-20 rule for eyes (every 20 minutes, look at something 20 feet away for 20 seconds).
  \item \textbf{Kindness and fairness.} We agree class rules together in week 1, give feedback as
        \enquote{two stars and a wish}, and give credit when we use someone else's ideas.
  \item \textbf{Accessible for everyone.} All materials use color-blind-safe colors, and there is a
        dyslexia-friendly edition with a clearer font and more spacing. Just ask.
\end{itemize}
\end{safety}

\section{What your child needs}
\begin{itemize}[itemsep=2pt]
  \item \textbf{Nothing to buy.} We provide computers, printed workbooks and stationery.
  \item \textbf{Homework (optional, 15--20 minutes a week).} Each set has a \enquote{Core} task and two
        optional extras. If you have a computer at home, Scratch and Thonny are free to install, and we can
        give you a copy of the course folder on a USB stick. If not, no problem: the paper parts can be done
        at the kitchen table, and computer tasks can be finished at the next session.
  \item \textbf{A good night's sleep and a snack before we start} help a lot more than you might think!
\end{itemize}

\section{How to help at home}
You do not need to know anything about coding to help. In fact, it often works best if you don't!
\begin{itemize}[itemsep=3pt]
  \item \textbf{Ask, don't fix.} Ask \enquote{What did you want it to do? What did it do instead?}
        Explaining a problem out loud often solves it.
  \item \textbf{Let your child be the teacher.} Ask them to show you one thing they made and how it works.
  \item \textbf{Celebrate bugs.} Mistakes in code are called bugs, and finding them is a real skill.
        Praise the strategy (\enquote{You tested it one step at a time, clever!}) more than the result.
  \item \textbf{Keep hands off the keyboard.} If you help, be the \enquote{navigator}: point and suggest,
        and let your child type.
  \item \textbf{Stop on a good note.} If frustration builds, take a break and come back later.
        Being stuck for a while is part of learning.
\end{itemize}

\section{Ask your child about\dots}
One question for each week. These make great conversation starters at dinner.

\begin{center}\small\renewcommand{\arraystretch}{1.45}
\begin{tabularx}{\linewidth}{@{}>{\bfseries}p{1.3cm}>{\RaggedRight}p{4.2cm}>{\RaggedRight\arraybackslash}X@{}}
\toprule
Week & Topic & Ask your child\dots\\
\midrule
1 & Robots Need Instructions & \enquote{How did you make the parent-robot draw a picture? Why does the order of steps matter?}\\
2 & Loop-de-Loop & \enquote{What is a loop? Can you show me a dance that uses a loop?}\\
3 & When This Happens\dots & \enquote{What does \emph{if} do in a program? Shall we play Simon-says-IF?}\\
4 & Keeping Score & \enquote{What is a variable? How did your game keep the score?}\\
5 & Messages and My Blocks & \enquote{What is a My Block? Why are color-blind-safe colors important?}\\
6 & Review Studio and Checkpoint & \enquote{Which bug was the hardest to find? How did you find it?}\\
7 & Hello, Python! & \enquote{What is different about typing code instead of using blocks? What does \enquote{Python is picky} mean?}\\
8 & Loops in Python & \enquote{Why do some lines in Python start with spaces?}\\
9 & Variables, Input and Choices & \enquote{What questions should a program never ask you?}\\
10 & Random Art and My Own Commands & \enquote{What is your final project going to be? What will it ask?}\\
11 & Project Studio & \enquote{What did your partner like about your project? What are you improving?}\\
12 & Showcase and Celebration & \enquote{What are you proudest of? What do you want to make next?}\\
\bottomrule
\end{tabularx}
\end{center}

\section{How to celebrate}
\begin{itemize}[itemsep=3pt]
  \item \textbf{Come to the showcase} in week 12. Each child runs their project for you, answers your
        questions, and explains a part of their code. You will get a feedback card to fill in
        (\enquote{two stars and a wish}).
  \item \textbf{Certificates.} Every child who completes Level 1 receives a printed certificate.
  \item \textbf{A home gallery.} Ask your child to describe their best picture and draw it, or
        print a screenshot of it, for the fridge. (Screenshots stay at home, not online.)
  \item \textbf{Celebrate effort, not just results.} \enquote{You kept trying until it worked} is the
        best thing a young coder can hear.
\end{itemize}

\bigskip
Thank you for trusting us with your child's learning. If you have any questions, or anything that
would help your child learn (for example, bigger text or more breaks), please tell me at pick-up time.

\bigskip
With warm wishes,\par\bigskip
\rule{6cm}{0.4pt}\par
{\small (Your coding club host)}

\vfill
\begin{center}\small\color{kcgray}
- - - - - - - - - - - - - - - - - - - - \faCut\ \ cut here and return (optional) - - - - - - - - - - - - - - - - - - - -
\end{center}
\noindent\textbf{Reply slip.}\quad Showcase: \tickbox\ we hope to come\quad \tickbox\ we cannot come, but say well done!\par\medskip
\noindent Something that helps my child learn: \rule{8.5cm}{0.4pt}\par\medskip
\noindent Questions: \rule{13cm}{0.4pt}
\end{document}
